\scanpage{109}
% archive sheet number: 108
% source 01-05, top
comme coordonnées affines et \deleted{\unsure{les}} exprimons
$\sigma_0,\sigma_1,\sigma_2$ p.r. à ces coordonnées --- pour
$\sigma_0,\sigma_2$ les coefficients sont des entiers bien déterminés
(les \unsure{signes|lignes} : $0$ ou $\pm1$) pour $\sigma_1$ il y a
des ``modules'' indéterminés $\lambda,\mu,\unsure{\nu|\gamma}$ satisfaisant
% source 06, upper middle
\[
\left\{\begin{aligned}
\nu^2&=1\\
\mu&=-\nu\lambda
\end{aligned}\right.\qquad\text{exprimant }\sigma_1^2=1
\]
% source 07, upper middle
puis $\nu=1$ (exprimant $u^3=\operatorname{id}$)
% source 08, upper right, corrected explanation
et enfin $\lambda^2+\lambda-1=0$
\deleted{(\unsure{car} $v=\sigma_1\sigma_2$, $v^{\unsure{5|6}}=\operatorname{id}$)}
\inserted{exprimant que $v^{\unsure{5|6}}=\operatorname{id}$.}
% source 09, middle, cancelled line
\deleted{On a fait \unsure{les calculs|le calcul} dans le cas $\nu=1$}
% source 10, middle, cancelled start
\deleted{$\mu=-\lambda$ (\unsure{qui est} \unsure{rappelé})}
% source 11-15, middle
--- au lieu de $\lambda$ on écrit aussi $\alpha'$ ... et on a trouvé,
exprimant $(\sigma_1\sigma_2)^5=\operatorname{id}$, l'équation
(en $\alpha$, ou $\alpha'$, au choix)
\[
\alpha^2+\alpha-1=0
\]
en termes de laquelle on a vérifié la fidélité.
% source 16-19, lower half, entire passage cancelled by large crosses
\deleted{Pour bien faire, il faudrait aussi traiter le cas où on ne suppose
pas $\nu=1$ i.e. \unsure{on n'impose pas|on ne supposerait pas}
$\sigma_1$ une réflexion (\unsure{en} car. rés. $\ne2$,
\unsure{cel.} revient à prendre $\nu=\deleted{\unsure{1|+1}}-1$)}
% source 20-22, cancelled continuation with local changes
\deleted{Considérons le \deleted{point} plan de $s_1,s_2,s_3$, soit $P$,
\inserted{d'équation $x+y+z=0$}
\deleted{\unsure{puisqu'}}\inserted{\unsure{notons qu'}} il contient
$s'_3=(\lambda,\mu,\nu)$ i.e. \deleted{\unsure{que} $\lambda+\mu+\nu=0$}
on a}
% source 23-26, bottom, all displays cancelled
\[
\deleted{\lambda+\mu+\nu=1}
\]
\deleted{i.e.}
\[
\deleted{\lambda-\nu\lambda+\nu=1\qquad\text{i.e. }\deleted{(\lambda-\nu\lambda)+\nu}}
\]
\[
\deleted{(\quad)\qquad(\lambda-1)(1-\nu)=0\qquad?}
\]
% source 27, bottom, locally crossed as well
\deleted{\deleted{Mais en fait, si $\nu\ne1$}}
